\documentclass{article}
\usepackage[UTF8]{ctex}
\usepackage{amsmath}
\usepackage{amssymb}

\begin{document}

\section{潜在赔付计算公式}

\subsection{1. 基本参数}
\begin{itemize}
    \item $S_0$：当前标的价格
    \item $K$：行权价
    \item $T$：到期时间（年）
    \item $\sigma$：隐含波动率
    \item $r$：无风险利率（通常为0）
\end{itemize}

\subsection{2. 标准差计算}
\[
\sigma_{\text{price}} = S_0 \times \sigma \times \sqrt{T}
\]

\subsection{3. 潜在赔付（看涨期权）}
对于看涨期权，潜在赔付为到期时标的价格超过行权价部分的期望值：
\[
E\left[\max(S_T - K, 0)\right] = \int_{-\infty}^{\infty} \max(S_T - K, 0) \cdot f(S_T) \, dS_T
\]
其中，$f(S_T)$ 是标的价格在到期时的概率密度函数（正态分布）：
\[
f(S_T) = \frac{1}{\sigma_{\text{price}} \sqrt{2\pi}} \exp\left(-\frac{(S_T - S_0)^2}{2\sigma_{\text{price}}^2}\right)
\]

\subsection{4. 潜在赔付（看跌期权）}
对于看跌期权，潜在赔付为行权价超过到期时标的价格部分的期望值：
\[
E\left[\max(K - S_T, 0)\right] = \int_{-\infty}^{\infty} \max(K - S_T, 0) \cdot f(S_T) \, dS_T
\]

\subsection{5. 数值积分实现}
通过数值积分近似计算，公式如下：

\subsubsection{看涨期权}
\[
\text{Potential Payout} = \sum_{i=1}^{n} \max(S_i - K, 0) \cdot f(S_i) \cdot \Delta S
\]

\subsubsection{看跌期权}
\[
\text{Potential Payout} = \sum_{i=1}^{n} \max(K - S_i, 0) \cdot f(S_i) \cdot \Delta S
\]

其中：
\begin{itemize}
    \item $n = 1000$（积分步数）
    \item $S_i = S_0 - 3\sigma_{\text{price}} + \frac{6\sigma_{\text{price}} \cdot i}{n}$（价格区间：$S_0 \pm 3\sigma_{\text{price}}$）
    \item $\Delta S = \frac{6\sigma_{\text{price}}}{n}$（步长）
\end{itemize}

\subsection{6. 完整公式}
\[
\text{Potential Payout} = 
\begin{cases} 
\sum_{i=1}^{n} \max(S_i - K, 0) \cdot \frac{1}{\sigma_{\text{price}} \sqrt{2\pi}} \exp\left(-\frac{(S_i - S_0)^2}{2\sigma_{\text{price}}^2}\right) \cdot \frac{6\sigma_{\text{price}}}{n} & \text{看涨期权} \\
\sum_{i=1}^{n} \max(K - S_i, 0) \cdot \frac{1}{\sigma_{\text{price}} \sqrt{2\pi}} \exp\left(-\frac{(S_i - S_0)^2}{2\sigma_{\text{price}}^2}\right) \cdot \frac{6\sigma_{\text{price}}}{n} & \text{看跌期权}
\end{cases}
\]

该公式表示：在期权被行权的情况下，期权卖方需要支付的平均内在价值。

\end{document}
